\subsection{平均不等式}

我们将对（有限或无限的）数值可测函数加强定理 1。

定义 1. --- 设 X 是一个局部紧空间，\(\mu\) 是 X 上的一个测度。给定一个在 X 中局部几乎处处有定义的数值函数 f（有限或无限），我们称函数 f 的依测度最大值或 \(\mu\)-最大值（相应地，依测度最小值或 \(\mu\)-最小值），记作 \(\mathrm{M}_{\infty}(f)\)（相应地，\(m_{\infty}(f)\)），指的是数 \(\alpha\)（使 \(f(x) \leqslant \alpha\)（相应地，\(f(x) \geqslant \alpha\)）局部几乎处处（关于 \(\mu\)）成立者）所成之集的下确界（相应地，上确界）。

由定义立得 \(m_{\infty}(f) = -M_{\infty}(-f)\)，于是由依测度最大值的每个性质可推出依测度最小值的一个相应性质。

对每个 \(\alpha > \mathrm{M}_{\infty}(f)\)，由 \(x \in \mathbf{X}\)（使 \(f(x) > \alpha\) 者）所成之集是局部可忽略的；而由 \(x \in \mathbf{X}\)（使 \(f(x) > \mathrm{M}_{\infty}(f)\) 者）所成之集是使 \(f(x) > r_n\) 的诸集之并（\(r_n\) 取遍有理数中 \(> \mathrm{M}_{\infty}(f)\) 者）；因此在局部几乎处处意义下 \(f(x) \leqslant \mathrm{M}_{\infty}(f)\)（§5, No.~2）。类似地，\(f(x) \geqslant m_{\infty}(f)\) 局部几乎处处成立；由此推出 \(m_{\infty}(f) \leqslant \mathrm{M}_{\infty}(f)\)（若测度 \(\mu\) 非零）；此外，关系式 \(m_{\infty}(f) = \mathrm{M}_{\infty}(f)\) 等价于 \(f\) 局部几乎处处等于一个常数。显然，若测度 \(\mu\) 非零，则

\[
\inf _ {x \in \mathrm{X}} f (x) \leqslant m _ {\infty} (f) \leqslant \mathrm{M} _ {\infty} (f) \leqslant \sup _ {x \in \mathrm{X}} f (x).
\]

若两个函数 \(f, g\) 局部几乎处处相等，则 \(m_{\infty}(f) = m_{\infty}(g)\)，且 \(\mathrm{M}_{\infty}(f) = \mathrm{M}_{\infty}(g)\) 。

最后，若 \(f\) 与 \(g\) 是使 \(f + g\) 局部几乎处处有定义的两个函数，则

\[
\mathrm{M} _ {\infty} (f + g) \leqslant \mathrm{M} _ {\infty} (f) + \mathrm{M} _ {\infty} (g)\tag{1}
\]

只要右端有定义，这由定义 1 立得；类似地，若 f 与 g 都 \(\geqslant0\) 且使 fg 局部几乎处处有定义，则

\[
\mathrm{M} _ {\infty} (f g) \leqslant \mathrm{M} _ {\infty} (f) \mathrm{M} _ {\infty} (g)\tag{2}
\]

只要右端有定义。

若 \(M_{\infty}(f) < +\infty\)，则 \(f(x) < +\infty\) 局部几乎处处成立，但未必几乎处处成立。称一个数值函数 f（对测度 \(\mu\) ）依测度有界，如果它几乎处处有定义且有限，而且数 \(m_{\infty}(f)\) 与 \(M_{\infty}(f)\) 都有限（后一条件相当于说 \(M_{\infty}(|f|)<+\infty\)）。

命题 1（平均不等式）. --- 设 f 是一个依测度有界的可测数值函数。对每个可积数值函数 \(g \geqslant 0\)，函数 fg（几乎处处有定义）可积，且

\[
m _ {\infty} (f) \int g d | \mu | \leqslant \int f g d | \mu | \leqslant \mathrm{M} _ {\infty} (f) \int g d | \mu |.\tag{3}
\]

此外，在不等式 (3) 的三个项中，除非在由 \(x \in X\)（满足 \(g(x) \neq 0\) 者）所成之集中，\(f\) 几乎处处等于 \(M_{\infty}(f)\) 或几乎处处等于 \(m_{\infty}(f)\)，否则其中两个项不可能相等。

事实上，\(fg\) 可测（§5, No.~3, 定理 1 的推论 5）；而且，不等式 \(m_{\infty}(f)g(x) \leqslant f(x)g(x) \leqslant M_{\infty}(f)g(x)\) 不仅局部几乎处处成立，甚至几乎处处成立，因为由点 \(x \in X\)（在那里 \(g(x) \neq 0\)）所成之集是可积集的可数并（§5, No.~6, 引理 1）。由此推出 \(fg\) 可积（§5, No.~6, 定理 5）且不等式 (3) 成立。另一方面，函数 \(M_{\infty}(f)g - fg\) 几乎处处有定义且等于 \((M_{\infty}(f) - f)g\)；因此它几乎处处 \(\geqslant 0\)（在 \(X\) 中）；由于关系式 \(M_{\infty}(f) \int g d|\mu| = \int fg d|\mu|\) 等价于 \(\int(M_{\infty}(f) - f)gd|\mu| = 0\)，故只有当函数 \((M_{\infty}(f) - f)g\) 可忽略时它才能成立，证毕。

撇开 \(\int g d|\mu| = 0\) 的平凡情形，不等式 (3) 可由应用于区间 \(D = [m_{\infty}(f), M_{\infty}(f)]\) 的 No.~1 定理 1 得出。可以对 No.~1 的定理 1 作出与命题 1 的那些补充类似的补充，它们刻画了点 \((\int f g d\mu)/(\int g d\mu)\) 属于 D 的边界的情形（习题 2）。

